% MathlibReview.tex -- status review of lean-DLMF from the perspective of a Mathlib
% reviewer, measured against the full review of 2026-10-30 (MathlibReview20261030.md).
% Compile with LuaLaTeX: latexmk -lualatex MathlibReview.tex
\documentclass[11pt,a4paper]{article}
\usepackage{leanreport}
\hypersetup{pdftitle={lean-DLMF: Mathlib-style status review, 2026-11-01}}

\newcommand{\dlmf}{\repofile{lean-DLMF}}

\begin{document}

\title{Mathlib-Style Status Review of \texttt{lean-DLMF}}
\author{Bastiaan J. Braams}
\date{2026-11-01}
\maketitle

\begin{abstract}
  This review checks the library at commit \texttt{9963e03} against the Mathlib reviewer
  questions: generality, redundant hypotheses, naming, docstrings, namespaces, module
  structure, proof length and reuse, file length and organization. It is a status update on the
  full review of 2026-10-30 (\dlmf{MathlibReview20261030.md}, at commit \texttt{3989f9e}). It
  measures what the review cleanup and the later proof rounds changed, records new findings,
  and lists the work to do before the next proof round. The earlier review's audit programs
  were run again unchanged on the current sources and compiled environment. The measurements
  and candidate lists are in \dlmf{MathlibReviewInventory20261101.json}. No library source was
  changed in this review.
\end{abstract}

\section{Summary}

The library builds (15{,}006 jobs). It has 1{,}210 library files with 545{,}049 lines and
26{,}344 explicit named declarations: 17{,}858 \texttt{theorem}, 5{,}012 \texttt{lemma},
3{,}352 \texttt{def}, 52 \texttt{abbrev}, 55 \texttt{structure}, 7 \texttt{inductive} and 8
named instances. There are 1{,}482 direct admissions, down from 1{,}553. In the compiled
environment, 1{,}581 theorems depend on \lean{sorryAx}, down from 1{,}655. No other
nonstandard axiom occurs.

These areas meet the Mathlib bar and need no further work:

\begin{itemize}
  \item Every library file is a \texttt{module} and has a module docstring before its first
    declaration. Every explicit declaration and all 205 directly declared structure fields
    have a formal docstring. (The utility \file{scripts/DumpDLMF.lean} is not a library module.)
  \item No file exceeds the 1{,}500-line limit, but one file is now exactly at it
    (\S\ref{sec:files}).
  \item The import survey finds no forbidden foundation-to-theorem imports, no missing modules,
    no file-import cycles and no new two-way area pairs; the 16 existing pairs are allowed.
  \item All 76 lowercase predicate definitions identified on 2026-10-30 were renamed, and no
    \texttt{Prop}-valued \texttt{def} now starts with a lowercase letter.
  \item The nine false formulations reported on 2026-10-30 have been corrected, and the two
    pending hypothesis questions (Leibniz rules, Haar's lemma) now carry explicit
    hypotheses.
\end{itemize}

These areas still need substantial work: development-round names (4{,}936 public
declarations), long proofs (913 over 60 lines), files near the length limit, stale module
headers, and the remaining duplicate and upstream-equivalent declarations. These are all
cosmetic or organizational. The review found no new false statements.

\section{Generality and hypotheses}

\paragraph{Closed hypotheses (new audit).}
This review adds one compiled check: a hypothesis whose type mentions none of the theorem's
variables is either provable, and therefore removable, or false, and therefore makes the
theorem vacuous. Across all project theorems there is exactly one such binder:
\begin{itemize}
  \item \lean{DLMF.Zeta.integral_quadratic_mul_cot_pi}
    (\dlmf{DLMF/Zeta/CotangentMoments.lean}, line~24) assumes
    $\lvert\pi x\cot(\pi x)\rvert\le1$ on $[0,\tfrac12]$. Its only caller, in
    \file{DLMF/Zeta/Sums.lean}, discharges the hypothesis with the existing lemma
    \lean{r25a_pi_mul_cot_bound}. Move that lemma down (or into a common Zeta support file)
    and prove the bound inside the theorem. The hypothesis is true, so the statement is not
    wrong, but it is not a finished interface.
\end{itemize}

\paragraph{Removable binders.}
The 2026-10-30 audit found 341 binders that occur in neither the type nor the elaborated
proof; the cleanup removed 224. The same audit now finds 125 binders in 109 admission-free
theorems (117 in public theorems). Of these, 117 are carried over from the earlier list and
were deliberately left for caller or API review. The 8 new ones are:
\begin{itemize}
  \item Four unused \lean{CompleteSpace} instances, in
    \lean{MeasureTheory.integral_comp_abs_eq_two_smul} (the neighbour of the binder removed
    from \lean{integral_eq_integral_Ioi_add_neg}),
    \lean{differentiableOn_integral_cexp_neg_mul_smul}, \lean{isBigO_integral_cexp_smul_sector}
    and \lean{Asymptotics.IsEquivalent.intervalIntegral_of_integral_norm}.
  \item An unused \lean{DecidableEq} instance in
    \lean{DLMF.EllipticIntegrals.contDiffAt_of_eqOn_analytic_comp_affine}.
  \item An unused \lean{NeZero} instance in
    \lean{DLMF.NumberTheory.DirichletCharacter.eq_mul_principal}.
  \item \lean{hloc : LocallyIntegrable φ} in
    \lean{DLMF.AlgebraicAnalytic.tendsto_gaussian_integral_of_growth}.
  \item An unused function argument \lean{ξ} in
    \lean{DLMF.AsymptoticApproximations.liouville_caseI}. An unused argument usually means the
    statement does not say what its docstring says, so check this one first.
\end{itemize}
Among the carried-over binders, those worth handling first are the ones that weaken
mathematically meaningful interfaces:
\begin{itemize}
  \item the slit-plane hypotheses on \lean{x + y} in the five Kummer addition theorems
    (\file{Confluent/ProductsAddition.lean});
  \item the two argument bounds and the non-integer condition on \lean{b} in
    \lean{wronskian_kummerU_kummer_pos} and \lean{wronskian_kummerU_kummer_neg};
  \item \lean{z ≠ 0}, \lean{z ≠ 1} in \lean{ordinaryHypergeometric_ode};
  \item \lean{cosh z ≠ 0} in four Jacobi modulus limits;
  \item the \lean{0 ≤ N} binders for natural~\lean{N} in the Airy remainder theorems.
\end{itemize}

\paragraph{Generality of the new shared APIs.}
The new \lean{ODE.wronskian} (normed-ring target over any nontrivially normed field) is
appropriately general. \lean{Function.mulDerivMul} and \lean{Function.oneSubMulDerivMul} are
restricted to $\mathbb{K}\to\mathbb{K}$. A normed-space target, as for the Wronskian, costs
nothing. The namespace \lean{Function} is also the wrong home for differential operators;
something like \lean{DiffOp} or the root namespace with a descriptive name would be better.
\lean{PowerSeries.HasSumAt} is still specific to $\ZZ$-coefficient series evaluated in
$\CC$. The 2026-10-30 suggestion to generalize it to a complete normed ring with a coefficient
homomorphism still stands. Its docstring says ``converges absolutely'' although the predicate
now stores only \lean{HasSum}. This is equivalent over $\CC$, but the docstring should say
which property is stored. Its declarations use explicit \lean{PowerSeries.}\ prefixes instead
of a \lean{namespace PowerSeries} block.

\section{Upstream reuse and duplication}

\paragraph{Exact upstream matches.}
58 project theorems have a type identical to a public theorem of Mathlib, TauCeti or
lean-codes, down from 62. These involve 33 distinct upstream declarations. The cleanup
removed the named examples, but the main family remains. There are 26 local copies of four
general facts:
\begin{itemize}
  \item 9 copies of \lean{Complex.ofReal_cpow_eq_exp}: in Bessel, Hypergeometric~(2), Legendre~(3),
    MatrixArgument, ErrorFunctions and the Fresnel transforms;
  \item 6 copies of \lean{Carlson.ofReal_cpow_half};
  \item 6 copies of \lean{Carlson.re_cpow_half_pos};
  \item 5 copies of \lean{Carlson.cpow_half_sq}.
\end{itemize}
There are also copies of \lean{Nat.stirlingSecond_succ_succ}, \lean{Nat.stirlingSecond_self},
\lean{Nat.ModEq.pow_totient}, \lean{Nat.Prime.of_mersenne}, \lean{Nat.totient_eq_mul_prod_factors},
\lean{Complex.besselJ_neg_int}, \lean{Complex.regularizedHGFun_zero},
\lean{Polynomial.derivative_bernoulli_add_one} and \lean{Finset.mem_inf}. Where a copy exists
only to carry a \texttt{@[dlmf]} tag, project policy calls for a \texttt{dlmf\_upstream}
pointer instead. The full list is in the inventory.

\paragraph{Identical definitions.}
17 pairs of public definitions have identical types and bodies, after excluding the
compiler-generated \lean{recOn}/\lean{casesOn} pairs. These are unambiguous duplicates to
consolidate:
\begin{itemize}
  \item \lean{Elementary.expPartialSum} and \lean{Confluent.expPartialSum};
  \item \lean{Combinatorics.partitionCountAtMostParts} and \lean{NumberTheory.partitionsAtMost};
  \item \lean{Confluent.pcfSect_logWedge} and \lean{Confluent.pcfLogWedge};
  \item \lean{Confluent.upperIncGammaOnLog} and \lean{Airy.airyUpperGammaOnLog};
  \item \lean{EllipticIntegrals.ellRedQuadU} and \lean{ei19b_redU};
  \item \lean{Zeta.q25b_slit} and \lean{z25b_slitOne};
  \item \lean{Zeta.z25a_gbinom} and \lean{AsymptoticApproximations.cbinom};
  \item \lean{AsymptoticApproximations.besselEnvCrit} and \lean{Legendre.besselEnvelopeCross};
  \item \lean{OrthogonalPolynomials.LegendreAddition.jet} and \lean{Legendre.legInt_D};
  \item \lean{Confluent.kzt_p} and \lean{OrthogonalPolynomials.laguerreWeight}.
\end{itemize}
The remaining pairs need a judgment call. They are coincident cut domains of different
inverse functions, coefficient families that coincide in one normalization, and
\lean{Selberg.vdmSq} versus \lean{siegelDiscriminant}. Those may keep separate descriptive
names, preferably with one defined as the other.

\section{Names and namespaces}

\paragraph{Development-round names.}
4{,}936 public declarations carry a round prefix (4{,}966 on 10-30), in 154 families. The
largest families are \lean{p22R1_} (237), \lean{p2p_} (150), \lean{p3q_} (133), \lean{p2m_}
(129), \lean{q2t_} (125), \lean{p19sR1_} (119) and \lean{p19lR1_} (117). By area: Mathieu
746, Combinatorics 453, JacobiElliptic 379, EllipticIntegrals 355, BernoulliEuler 337, Zeta
293, Lam\'e 291, CoalescingSaddles 280. Thirteen type or predicate names carry the same
prefixes (for example \lean{P2lIsSol}, \lean{Q25bIsLerchCont}, \lean{P3jDich},
\lean{Q2mCoeffCont}, \lean{S26aGood}, \lean{Z25bLerchParamOK}). 59 file paths contain round
identifiers (60 before). This is the largest gap between the library and Mathlib style. Most
of these declarations are helpers, so the cheapest fix is often to make them
\texttt{private} rather than rename them. The library now has 25{,}744 public and only 967
private declarations. Renaming should be done area by area, together with the visibility
pass.

\paragraph{Theorem names in camel case.}
42 public theorems have a lowerCamelCase name with no underscore. About half are correct,
because they name a predicate in a conclusion (for example \lean{IsPIIOn.hasDerivAt}). The
others name a theorem after a person or concept rather than its conclusion. Mathlib would
ask for snake case describing the statement:
\lean{AlgebraicAnalytic.mittagLeffler}, \lean{poissonSummation}, \lean{weierstrassProduct},
\lean{lagrangeInversion}, \lean{lagrangeBurmann},
\lean{IsSimpleClosedContour.residueTheorem}, \lean{argumentPrinciple},
\lean{cauchyIntegralFormula}, \lean{QHypergeometric.generalBaileyTransform},
\lean{limitingBaileyTransform}, \lean{rogersBaileyPair}, \lean{Bailey.watsonR},
\lean{Struve.trigBeta}, \lean{RamanujanAux.gammaShiftBound},
\lean{HasAsymptoticPowerSeries.mcMahon}, \lean{IsPainleveVIOn.ellipticForm},
\lean{IsPIIOn.shiftUp}/\lean{shiftDown}, \lean{IsPainleveIIOn.tPlus}/\lean{tMinus} and
\lean{IsLinODE2SolOn.linComb}.

\paragraph{Uppercase data definitions.}
18 public data definitions have UpperCamelCase names. All are one-letter or abbreviated
helper quantities inside auxiliary namespaces: \lean{DixonAux.H}, \lean{K}, \lean{P},
\lean{Lh}, \lean{Rh}; \lean{QGammaLimit.D}, \lean{K}, \lean{L}; \lean{QGammaMono.G},
\lean{T}; \lean{JacobiDisc.EE}, \lean{SS}, \lean{Fr}; \lean{LagDisc.SSL}, \lean{KL};
\lean{HerDisc.SSH}; \lean{WaveEquation.F01}; \lean{SphericalHarmonics.SU2}. They should be
\texttt{private} or receive lowerCamel descriptive names. The \lean{…Aux} namespaces
themselves would not pass Mathlib review.

\paragraph{Remaining namespace items from 10-30.}
The root-level \lean{stRem_*}, \lean{r27_*} and \lean{q26b_*} APIs in \file{DLMFMathlib}
(40 declarations) still need mathematical namespaces and names. The cleanup fixed the
nested \lean{DLMF.NumberTheory.DirichletCharacter} namespace. The Airy remainder family
still uses \lean{airy_ai_…}/\lean{airy_bi_…} (20 theorems) instead of the
\lean{airyAi}/\lean{airyBi} vocabulary.

\section{Documentation}

Formal coverage is complete (see the Summary). Semantic quality has three outstanding
issues:
\begin{itemize}
  \item \textbf{Stale acquisition headers.} 85 module docstrings still contain acquisition
    boilerplate (``Statements acquired from the NIST DLMF, checked \ldots'', ``New statements
    are initially admitted''). By area: Legendre 24, Airy 13, ParabolicCylinder 11,
    IncompleteGamma 11, Struve 8, ErrorFunctions 8. Twenty-three of these files contain no
    \texttt{sorry} at all. Example: \file{DLMF/Airy/Generalized.lean} (1{,}342 lines) has a
    three-sentence header that names no result. Replace these headers with a statement of
    scope, conventions and main results.
  \item \textbf{Short headers on large files.} 243 files over 500 lines have module
    docstrings under 50 words. Not every short header is inadequate, but Mathlib expects a
    ``Main results'' or ``Main definitions'' list in a file of this size. Many good examples
    already exist, such as \file{MatrixArgument/GammaBeta.lean} and \file{Zeta/Sums.lean}.
  \item \textbf{Small errors.} \file{GeneralizedHypergeometric/MathlibCoverage.lean} has
    ``Chapter16'' in its title. The \lean{PowerSeries.HasSumAt} docstring is noted above.
\end{itemize}
The 10-30 items on \file{Core/Framework.lean}, \file{DLMF.lean},
\file{Painleve/Basic/Equations.lean} and \file{Hypergeometric/ConnectionFormulas.lean} have
been fixed.

\section{Proofs: length, repetition and reuse}

There are 913 admission-free tactic proofs over 60 lines (901 on 10-30) and 214 over 100
(208). The 14 new proofs over 60 lines come from the latest proof rounds. The largest are
\lean{symEllIntRD_asymp_small_x} (137 lines), \lean{symEllIntRG_asymp_small_x} (131),
\lean{integral_quadratic_mul_cot_pi} (127), \lean{integral_symEllIntRD_interp} (109),
\lean{qLaguerre_monomial_orthogonality} (105) and
\lean{hasSum_riemannZeta_two_mul_div_odd_even} (103). The fragment scan pairs the
$R_G$/$R_D$ small-$x$ asymptotics (four shared eight-line windows), the two odd-zeta sums in
\file{Zeta/Sums.lean}, and \lean{symEllIntRG_add_symEllIntRG} with
\lean{symEllIntRG_duplication}. The $R_F$, $R_G$ and $R_D$ small-$x$ proofs should share one
lemma on the logarithmic singularity of the defining integral.

The refactoring targets of 10-30 are unchanged. Each remains the most valuable proof cleanup
in its area:
\begin{enumerate}
  \item The logarithmic connection formulas: \lean{logarithmic_connection_one_sub_recip}
    (268 lines), \lean{…_recip_recipGamma} (266), \lean{…_recipGamma} (239) and
    \lean{logarithmic_connection_one_sub} (236).
  \item The Jacobi cyclic certificates: \lean{jacobi_cyclic_scdd_three} (270),
    \lean{jacobi_cyclic_sdcc_three} (257), \lean{jacobi_cyclic_d_pairs} (260) and
    \lean{jacobi_cyclic_d2_triple} (200).
  \item The Painlev\'e~VI B\"acklund transformations: \lean{IsPainleveVIOn.quartic_root} (260)
    and \lean{quadratic} (197).
  \item The two Liouville--Green Volterra constructions (211 and 208 lines).
\end{enumerate}
The Mathieu and Jacobi-quotient repetitions were reduced in the cleanup. The fragment scan
still finds 489 repeated eight-line windows (495 before). The most affected files are
\file{Mathieu/SmallQCeSe.lean} (20), \file{Legendre/SpecialValues.lean} (15),
\file{Mathieu/SmallQGeneral.lean} (13), \file{JacobiElliptic/Addition/ThetaAndQuotients.lean}
(11) and \file{Confluent/ParabolicCylinder/WronskiansAndConnections.lean} (10).

The successful build emits 1{,}310 warnings unrelated to admissions (1{,}468 before). Of
these, 618 are unused \texttt{simp} arguments, 303 are unused-tactic reports, 132 are
never-executed tactics, 84 are unused variables, 72 are ``used \texttt{X} where \texttt{Y}
would suffice'' and 60 are deprecations. The deprecations should be fixed before the next
Mathlib bump.

\section{Files and organization}\label{sec:files}

138 library files exceed 1{,}000 lines (134 before). Several grew past the point where the
next proof round will break the limit:
\begin{center}
\begin{tabular}{lrr}
  \toprule
  File & 10-30 & now \\
  \midrule
  \file{MatrixArgument/GammaBeta.lean} & 1{,}454 & \textbf{1{,}500} \\
  \file{Zeta/Sums.lean} & 1{,}372 & 1{,}474 \\
  \file{QHypergeometric/BasicSummations.lean} & 1{,}412 & 1{,}463 \\
  \file{QHypergeometric/Calculus.lean} & 1{,}331 & 1{,}427 \\
  \file{Struve/Lommel.lean} & 1{,}312 & 1{,}382 \\
  \file{QHypergeometric/Convergence.lean} & 1{,}257 & 1{,}351 \\
  \file{EllipticIntegrals/SymmetricAddition.lean} & --- & 1{,}351 \\
  \bottomrule
\end{tabular}
\end{center}
\file{Spheroidal/Angular.lean} (1{,}450), \file{Hypergeometric/ConnectionFormulas.lean}
(1{,}446), \file{Confluent/LargeA.lean} (1{,}438), \file{Legendre/AdditionTools.lean}
(1{,}419), \file{Elementary/Trigonometric/InverseIdentities.lean} (1{,}417) and
\file{Legendre/Series.lean} (1{,}406) are also near the limit. The splits proposed on 10-30
for \file{GammaBeta} (Laplace transform, gamma integrals, beta integrals) and for
\file{Spheroidal/Angular} are now urgent. \file{Zeta/Sums.lean} separates naturally into the
power series 25.8.5--25.8.8 and the even-zeta summations; the new
\file{CotangentMoments.lean} belongs with the latter. The four Q-hypergeometric files would
benefit from a shared lower-level file for the shift and derivative machinery, as proposed
on 10-30.

Import hygiene is unchanged: 430 files have an unsorted consecutive import block, 57 files
import \lean{Mathlib.Tactic.*} modules directly (352 lines), and 4{,}470 source lines exceed
100 characters. Runs of four or more blank lines are gone. The 10-30 experiment showed that
the 18 tactic imports in \file{SegmentPrimitive.lean} were removable, and they have been
removed. \file{Painleve/Basic/Equations.lean} still has 16.

\paragraph{Follow-up.} Later on 2026-11-01 the four files of item~1 below were split along their
mathematical sections, with the original modules kept as umbrellas; the largest piece has 707
lines. See \dlmf{REVIEW.md}, ``Splitting the files at the length limit''.

\section{Recommended order before the next proof round}

\begin{enumerate}
  \item \textbf{Split the files at or near the limit:} \file{MatrixArgument/GammaBeta.lean}
    first, then \file{Zeta/Sums.lean}, \file{QHypergeometric/BasicSummations.lean} and
    \file{QHypergeometric/Calculus.lean}. Otherwise proof work in those areas will
    immediately break the 1{,}500-line rule.
  \item \textbf{Fix the small verified interface defects:} the closed hypothesis in
    \lean{integral_quadratic_mul_cot_pi}, the 8 new removable binders (check
    \lean{liouville_caseI} first), and the carried-over binders on the Kummer addition and
    Wronskian theorems.
  \item \textbf{Replace the 26 copies of the four complex-power facts} and the other exact
    upstream copies with the upstream declarations or \texttt{dlmf\_upstream} pointers.
    Consolidate the ten clear duplicate definitions.
  \item \textbf{Run a combined visibility and naming pass, area by area,} starting with
    Mathieu, Combinatorics, JacobiElliptic and EllipticIntegrals. Make one-use helpers
    private, rename the reusable ones, and fix the camel-case theorem names and the
    uppercase helper definitions.
  \item \textbf{Replace the 85 stale acquisition headers} with scope and main-results summaries.
  \item The long-proof refactorings and the linter and import cleanup can run alongside proof
    work. They do not block it.
\end{enumerate}

\section*{Method and validation}

The ordinary build command \texttt{lake build} completed (15{,}006 jobs). The 2026-10-30 audit
programs \file{scan\_source.py} and \file{CompiledReview.lean}, stored in the earlier
inventory, were rerun unchanged except for output paths and the file listing. The compiled
audit read 27{,}069 compiled constants. A docstring probe confirmed that the
apparent ``missing'' structure docstrings in its raw output are constructor rows. The new
closed-hypothesis audit (\file{ClosedHyps.lean}, stored in the new inventory) scans all
project theorems in the compiled environment. \file{scripts/import\_survey.py} and
\file{scripts/check\_file\_lengths.py} pass. The removable-binder findings rely on the binder
being absent from the elaborated type and proof term. Unlike on 10-30, this review did not
kernel-check each narrowed declaration individually. Counts of naming, fragment and
documentation candidates come from heuristics and are not adjudicated violations.

\end{document}
